% P17-D2: monetary risk, comparisons, and wealth invariance.
\chapter{Monetary Risk and Risk Aversion}
\label{sec:p17-risk-aversion}

Expected utility permits precise questions about monetary risk. Which
lottery is preferred by every increasing utility? When is a sure mean
preferred to a lottery? What does it mean for one preference to be more
risk averse than another? This chapter develops Rubinstein's monetary
lottery, certainty-equivalent and comparative-risk treatments
\cite[2015 revision, pp.~108--118]{Rubinstein2015LectureNotes}.
All stated mathematical results have complete proofs in the annex.

\section{The monetary domain and stochastic dominance}
\label{sec:p17-monetary-lotteries}

\begin{hbdefinition}{def:p17-monetary-lottery}{Finite-support lotteries on an interval}
Let \(I\subseteq\mathbb R\) be a nonempty interval of monetary consequences.
Write \(\Delta_f(I)\) for the functions \(p:I\to[0,1]\) with finite
support \(\operatorname{supp}(p)=\{x:p(x)>0\}\) and total mass one.
All sums below range over finite supports, including unions of supports
when two lotteries are compared. Put
\[
       \mu(p)\defeq\sum_xp(x)x,\qquad
       U_u(p)\defeq\sum_xp(x)u(x)\quad(u:I\to\mathbb R).
\]
Since an interval contains finite convex combinations of its points,
\(\mu(p)\in I\). The certain lottery at \(x\in I\) is \(\delta_x\).
An expected-utility preference \(\succeq_u\) is defined by
\(p\succeq_u q\) exactly when \(U_u(p)\geq U_u(q)\).

This chapter assumes a specified expected-utility function on the whole
monetary interval, as in the source's opening scope
\cite[p.~108]{Rubinstein2015LectureNotes}. The finite-prize representation
theorem does not, by itself, prove the existence of one consistent function
on every finite subset of an infinite interval. Here \(u\) is supplied;
its regularity assumptions are stated for each result.

\emph{Strict monotonicity in money} means
\(\delta_x\succ_u\delta_y\) whenever \(x>y\). For an expected-utility
preference this is exactly strict increase of \(u\).
If \(I\) is a singleton, its only lottery is certain; all subsequent
mean and certainty comparisons are trivial on that domain.
\index{lottery!monetary}\index{risk aversion!monetary domain}
\end{hbdefinition}

\begin{hbdefinition}{def:p17-first-order-dominance}{First-order stochastic dominance}
For \(p,q\in\Delta_f(I)\), write \(p\succeq_{\rm FSD}q\) when
\(U_u(p)\geq U_u(q)\) for every continuous strictly increasing
\(u:I\to\mathbb R\). Define, for every real threshold \(t\),
\[
        F_p(t)\defeq\sum_{x\leq t}p(x),\qquad
        G_p(t)\defeq\sum_{x\geq t}p(x).
\]
These are respectively the cumulative and upper-tail probabilities,
with their indicated weak endpoint inequalities. At an atom,
\(G_p(t)\) is not generally \(1-F_p(t)\).
The source denotes this dominance relation by \(D_1\)
\cite[pp.~109--110]{Rubinstein2015LectureNotes}.
\index{stochastic dominance!first order}
\end{hbdefinition}

\begin{hbtheorem}{thm:p17-finite-stochastic-dominance}{proved=annex:p17.finite-stochastic-dominance}{The finite-support dominance criterion}
For finite-support lotteries on an interval, the following are equivalent:
\begin{enumerate}
\item \(p\succeq_{\rm FSD}q\).
\item \(G_p(t)\geq G_q(t)\) at every real \(t\).
\item \(F_p(t)\leq F_q(t)\) at every real \(t\).
\end{enumerate}
If the distinct points in the union of the supports are
\(x_0<\cdots<x_K\), then for any real function on those points,
\[
 U_u(p)=u(x_0)+\sum_{k=1}^K
              G_p(x_k)\bigl(u(x_k)-u(x_{k-1})\bigr).
\]
If the tail criterion fails, a continuous strictly increasing piecewise
linear utility witnessing \(U_u(p)<U_u(q)\) can be constructed explicitly.
Consequently the same dominance criterion also applies if all
nondecreasing utilities are used instead.
\end{hbtheorem}

\begin{hbexample}{def:p17-dominance-and-means}{Dominance and equal means}
On \(I=\mathbb R\), let \(p=\tfrac12\delta_1+\tfrac12\delta_3\) and
\(q=\tfrac12\delta_0+\tfrac12\delta_2\). Each prize of \(q\) has
been shifted up by one, so every strictly increasing utility gives a larger
expectation under \(p\). The tail and cumulative criteria confirm this
comparison at the four support points.

By contrast, \(r=\tfrac12\delta_0+\tfrac12\delta_2\) and
\(s=\delta_1\) both have mean one, and neither dominates the other.
For \(0\leq t<1\), \(F_r(t)=1/2>0=F_s(t)\); for
\(1\leq t<2\), \(F_r(t)=1/2<1=F_s(t)\). Their comparison requires
additional information about the utility's curvature. A mean alone
does not determine a lottery preference.
\end{hbexample}

\section{Concavity and aversion to a mean-preserving risk}
\label{sec:p17-risk-concavity}

\begin{hbdefinition}{def:p17-concavity-risk}{Concavity and risk aversion}
A function \(u:I\to\mathbb R\) is \emph{concave} if
\[
 u(\alpha x+(1-\alpha)y)
       \geq\alpha u(x)+(1-\alpha)u(y)
          \quad(x,y\in I,\ 0\leq\alpha\leq1).
\]
It is \emph{strictly concave} if the inequality is strict for
\(x\ne y\) and \(0<\alpha<1\).
The preference \(\succeq_u\) is \emph{risk averse} if
\(\delta_{\mu(p)}\succeq_u p\) for every \(p\in\Delta_f(I)\),
and \emph{risk neutral} if all these comparisons are indifferences.
These are weak risk aversion and exact mean-indifference, respectively;
strict preference for every nonconstant lottery's sure mean is a
stronger property. The source's risk-aversion definition uses the weak
comparison \cite[pp.~110--111]{Rubinstein2015LectureNotes}.
\end{hbdefinition}

\begin{hblemma}{thm:p17-finite-jensen}{proved=annex:p17.finite-jensen}{Finite Jensen inequality and secant slopes}
If \(u\) is concave, \(\alpha_k\geq0\), \(\sum_{k=1}^n\alpha_k=1\),
\(n\geq1\), and \(x_k\in I\), then
\[
             u\left(\sum_k\alpha_kx_k\right)
                        \geq\sum_k\alpha_ku(x_k).
\]
For strictly concave \(u\), the inequality is strict if at least two
distinct \(x_k\) have positive weight. For any \(a<b<c\) in \(I\),
concavity gives
\[
 \frac{u(c)-u(b)}{c-b}
 \leq\frac{u(c)-u(a)}{c-a}
 \leq\frac{u(b)-u(a)}{b-a}.
\]
On an open interval, a twice continuously differentiable function is
concave exactly when \(u''\leq0\) everywhere. Continuity at an
included endpoint is an additional assumption when needed below;
the open-interval derivative criterion makes no assertion about an
uncontrolled endpoint value.
\end{hblemma}

\begin{hbtheorem}{thm:p17-risk-aversion-concavity}{proved=annex:p17.risk-aversion-concavity}{Expected-utility risk aversion and concavity}
For \(u:I\to\mathbb R\), the preference \(\succeq_u\) is risk
averse if and only if \(u\) is concave. Strict concavity makes the
sure mean strictly preferred for every lottery with at least two
different prizes of positive probability.
The preference is risk neutral if and only if \(u\) is affine on
\(I\); when it is also strictly increasing and \(I\) is nondegenerate,
its affine slope is positive.
\end{hbtheorem}

\section{Certainty equivalents and comparison of risk attitudes}
\label{sec:p17-certainty-equivalents}

\begin{hbdefinition}{def:p17-certainty-equivalent}{Certainty equivalent and risk premium}
For a lottery \(p\), a \emph{certainty equivalent} is a consequence
\(c\in I\) satisfying \(\delta_c\sim_u p\). When it exists uniquely,
write it \(\operatorname{CE}_u(p)\). Its \emph{risk premium} is
\[
                \pi_u(p)\defeq\mu(p)-\operatorname{CE}_u(p).
\]
These quantities have units of the monetary consequence, rather than
the units of a chosen utility scale
\cite[p.~111]{Rubinstein2015LectureNotes}.
\index{certainty equivalent}\index{risk premium}
\end{hbdefinition}

\begin{hbtheorem}{thm:p17-certainty-equivalent}{proved=annex:p17.certainty-equivalent}{Existence, uniqueness and the sign of the premium}
If \(u\) is continuous and strictly increasing on \(I\), every
finite-support lottery has a unique certainty equivalent, with
\[
 \operatorname{CE}_u(p)=u^{-1}(U_u(p)),\qquad
 \min\operatorname{supp}(p)\leq\operatorname{CE}_u(p)
                             \leq\max\operatorname{supp}(p).
\]
Here \(u^{-1}\) is defined on the interval \(u(I)\), and the
expected utility always lies in that interval. The condition
\(\pi_u(p)\geq0\) for every lottery is equivalent to concavity of \(u\). Under strict
concavity, \(\pi_u(p)>0\) for each nonconstant lottery. Positive
affine changes of \(u\) preserve both \(\operatorname{CE}_u\) and
\(\pi_u\).
\end{hbtheorem}

\begin{hbdefinition}{def:p17-more-risk-averse}{More risk averse on a common monetary domain}
Let \(u_1,u_2\) be continuous strictly increasing functions on the
same interval \(I\). Say that preference 1 is \emph{at least as risk
averse as} preference 2 if, for every \(p\in\Delta_f(I)\) and every
\(c\in I\),
\[
                  p\succeq_{u_1}\delta_c
                         \quad\Longrightarrow\quad
                  p\succeq_{u_2}\delta_c.
\]
Thus every lottery acceptable to preference 1 against a stated sure
consequence is also acceptable to preference 2. Neither preference is
required to be risk averse relative to the mean merely for this
comparative definition to apply
\cite[pp.~111--113]{Rubinstein2015LectureNotes}.
\index{risk aversion!comparative}
\end{hbdefinition}

\begin{hbtheorem}{thm:p17-comparative-risk}{proved=annex:p17.comparative-risk}{Three equivalent comparisons of risk aversion}
For the two utilities just specified, the following are equivalent:
\begin{enumerate}
\item Preference 1 is at least as risk averse as preference 2.
\item \(\operatorname{CE}_{u_1}(p)\leq\operatorname{CE}_{u_2}(p)\)
for every \(p\in\Delta_f(I)\).
\item The function \(\phi=u_1\circ u_2^{-1}\) is concave on
the interval \(u_2(I)\).
\end{enumerate}
Equivalently, every lottery has \(\pi_{u_1}(p)\geq\pi_{u_2}(p)\).
The comparison includes ties and needs no differentiability.
\end{hbtheorem}

\begin{hbexample}{def:p17-certainty-example}{Two utilities for the same lottery}
On \(I=(0,\infty)\), let \(u_1(x)=\log x\) and \(u_2(x)=\sqrt x\).
Their relative transformation is \(\phi(t)=2\log t\) for \(t>0\),
which is concave. Thus preference 1 is at least as risk averse as
preference 2. For the hypothetical lottery
\(p=\tfrac12\delta_1+\tfrac12\delta_9\),
\[
\begin{aligned}
 \mu(p)&=5,&
 \operatorname{CE}_{u_1}(p)&=\exp\bigl(\tfrac12\log1+\tfrac12\log9\bigr)=3,\\
 \operatorname{CE}_{u_2}(p)&=\bigl(\tfrac12\sqrt1+\tfrac12\sqrt9\bigr)^2=4,&
 \pi_{u_1}(p)&=2,\qquad \pi_{u_2}(p)=1.
\end{aligned}
\]
The numbers follow from the stipulated utility functions. They do not
estimate an individual's preferences or establish that one attitude is
normatively preferable.
\end{hbexample}

\section{The differentiable coefficient and local indifference curves}
\label{sec:p17-absolute-risk-aversion}

\begin{hbdefinition}{def:p17-arrow-pratt}{Coefficient of absolute risk aversion}
Let \(I\) be an open interval, \(u\) twice continuously differentiable,
and \(u'(x)>0\) at every point. Define
\[
                             r_u(x)\defeq-\frac{u''(x)}{u'(x)}.
\]
The positivity condition makes the quotient defined. Strict increase
alone does not imply an everywhere positive derivative. This is the
absolute risk-aversion coefficient used in the source's differentiable
comparison \cite[p.~114]{Rubinstein2015LectureNotes}.
\index{risk aversion!Arrow--Pratt coefficient}
\end{hbdefinition}

\begin{hbtheorem}{thm:p17-arrow-pratt-comparison}{proved=annex:p17.arrow-pratt-comparison}{Differential comparison and local curvature}
For \(u_1,u_2\) satisfying the preceding differentiability and positivity
conditions on the same open interval, preference 1 is at least as risk
averse as preference 2 exactly when
\(r_{u_1}(x)\geq r_{u_2}(x)\) everywhere. The coefficient is unchanged
by positive affine changes of utility.

For one such \(u\), fix \(0<\theta<1\) and \(t\in I\). The
indifference curve through the certain outcome \((t,t)\) for the
two-prize lottery \(\theta\delta_x+(1-\theta)\delta_y\) is locally
the twice continuously differentiable function \(y=\psi_t(x)\), where
\[
 \theta u(x)+(1-\theta)u(\psi_t(x))=u(t).
\]
It satisfies
\[
 \psi_t(t)=t,\qquad
 \psi_t'(t)=-\frac{\theta}{1-\theta},\qquad
 \psi_t''(t)=r_u(t)\frac{\theta}{(1-\theta)^2}.
\]
These are local derivatives; a comparison of two coefficients at one
wealth level alone does not establish the global comparative relation.
\end{hbtheorem}

For the preceding logarithmic and square-root example,
\(r_{u_1}(x)=1/x\) and \(r_{u_2}(x)=1/(2x)\). Their pointwise
ordering agrees with the already obtained certainty-equivalent comparison.

\section{Wealth changes and translation invariance}
\label{sec:p17-wealth-invariance}

\begin{hbdefinition}{def:p17-wealth-translation}{Translation of a lottery and invariance to wealth}
For this section take \(I=\mathbb R\). If \(p\in\Delta_f(\mathbb R)\)
and \(w\in\mathbb R\), let \(T_wp\) be its translated lottery:
\((T_wp)(y)=p(y-w)\). A preference is \emph{invariant to wealth}
when
\[
            p\succeq q\quad\Longleftrightarrow\quad
                 T_wp\succeq T_wq
          \qquad(p,q\in\Delta_f(\mathbb R),\ w\in\mathbb R).
\]
The domain \(\mathbb R\) makes every translation admissible. If prizes
are interpreted as changes from initial wealth \(w\), deriving their
comparison from a single preference over final wealth uses
\(U_u(T_wp)=\sum_xp(x)u(w+x)\). This interpretation and the further
invariance condition are distinct assumptions in the source
\cite[pp.~115--117]{Rubinstein2015LectureNotes}.
\end{hbdefinition}

\begin{hbtheorem}{thm:p17-cara-characterization}{proved=annex:p17.cara-characterization}{Continuous wealth-invariant risk aversion}
Let \(u:\mathbb R\to\mathbb R\) be continuous and strictly increasing.
Its expected-utility preference is risk averse and invariant to wealth
if and only if \(u\) has one of the forms
\[
               u(x)=a+b x\quad(b>0),
       \qquad\text{or}\qquad
               u(x)=a-b e^{-\rho x}\quad(b>0,\ \rho>0).
\]
The first has \(r_u=0\); the second has \(r_u=\rho\). Thus
differentiability follows from the conclusion and need not be added
to the assumptions of this characterization. These preferences are
called constant absolute risk aversion, including the risk-neutral case.
For either form,
\(\operatorname{CE}_u(T_wp)=w+\operatorname{CE}_u(p)\), and the
risk premium is unchanged by the wealth translation.
\end{hbtheorem}

The source gives a grid argument and a separate differentiable argument
for this characterization \cite[pp.~117--118]{Rubinstein2015LectureNotes}.
The annex retains those two routes and supplies the full continuous-domain
functional-equation proof, including extension of affine uniqueness to
finite-support lotteries on the whole line. It explicitly justifies
compatibility across grids and the passage to all real wealth levels.

\begin{hbexample}{def:p17-cara-premium}{A symmetric risk under exponential utility}
For \(u(x)=-e^{-\rho x}\), \(\rho>0\), consider
\(p=\tfrac12\delta_{w-d}+\tfrac12\delta_{w+d}\), with \(d\geq0\).
Define \(\cosh t=(e^t+e^{-t})/2\). Then direct averaging and inversion give
\[
 U_u(p)=-e^{-\rho w}\cosh(\rho d),\qquad
 \operatorname{CE}_u(p)=w-\frac{\log\cosh(\rho d)}{\rho},\qquad
 \pi_u(p)=\frac{\log\cosh(\rho d)}{\rho}.
\]
The premium is independent of \(w\), is zero for \(d=0\), and is
positive for \(d>0\), since \(\cosh t>1\) for \(t\ne0\).
For \(\rho=\log2\) and \(d=1\) it is
\(\log(5/4)/\log2\). This exact finite calculation does not involve
an approximation for small risks.
\end{hbexample}

\begin{hbwarning}{def:p17-risk-domain-boundaries}{Consequences and the scope of a risk comparison}
Choosing final wealth as the consequence identifies situations that
produce the same final-wealth lottery. Starting at 2000 with a sure
loss of 500, or at 1000 with a sure gain of 500, gives the same certain
1500. A fifty-fifty loss of 1000 or zero from the first starting point,
and a fifty-fifty gain of 1000 or zero from the second, both give the
same final lottery on 1000 and 2000. A single preference over final
wealth therefore makes the same comparison in these two problems.
Treating gains and losses as different consequence descriptions changes
that premise \cite[pp.~115--116]{Rubinstein2015LectureNotes}.

All lotteries in this chapter have finite support. No St.~Petersburg
limit, infinite stopping procedure, general continuous distribution,
or empirical calibration claim follows merely from the finite-sum
identities. The source's separate discussion of Rabin's calibration
argument at pp.~118--119 requires its own full statement and proof.
Higher-order stochastic dominance, ambiguity, and risk sensitivity to
information likewise remain separate treatments.
\end{hbwarning}
